\section{Casos de fallo y manual de depuración: convertir un sistema inutilizable en un sistema iterable}

\subsection{Estratificación de tipos de fallo: primero agrupar, luego corregir}
En los proyectos de Agentes multimodales, el fallo no es la excepción sino la norma. Una idea de ingeniería clave implícita en el curso es: no fijarse directamente en el «fallo final», sino primero descomponer el fallo en subcategorías corregibles. Solo cuando los tipos de fallo se agrupan en categorías, las acciones de optimización resultan eficaces.

\begin{importantbox}{Clasificación recomendada de fallos en cuatro niveles}
\begin{enumerate}
    \item \textbf{Fallo de percepción}: no entender la interfaz, no leer el texto clave, ubicar incorrectamente el objeto objetivo;
    \item \textbf{Fallo de planeación}: orden incorrecto de subobjetivos, pasos omitidos, pasos repetidos en exceso;
    \item \textbf{Fallo de ejecución}: el formato de la acción es correcto pero no coincide con el contexto, o falla la llamada a la interfaz;
    \item \textbf{Fallo de verificación}: la tarea está casi completa, pero el script de verificación no coincide con la definición del objetivo.
\end{enumerate}
\end{importantbox}

\begin{knowledgebox}{Por qué «agrupar los fallos en categorías» afecta directamente la eficiencia del entrenamiento}
Si todas las muestras de fallo se reincorporan de manera uniforme, se amplifica el ruido. Una vez agrupadas por categorías, los datos pueden reforzarse de manera dirigida según el tipo de problema: para los fallos de percepción, muestras de grounding y OCR; para los fallos de planeación, muestras de razonamiento en cadenas largas; para los fallos de ejecución, muestras de restricción de acciones; para los fallos de verificación, muestras de prueba de robustez del script.
\end{knowledgebox}

\subsection{Ciclo cerrado de depuración: registros, reproducción y experimentos de control}
La práctica de ingeniería correspondiente al contenido del curso puede resumirse en tres pasos: primero, capturar el registro completo; segundo, reproducir la trayectoria; tercero, realizar un experimento de control mínimo. El objetivo del experimento de control no es obtener la mejor puntuación, sino confirmar si un cambio determinado realmente reduce cierto tipo de fallo.

\begin{table}[H]
\centering
\begin{tabular}{p{2.6cm}p{4.6cm}p{4.0cm}}
\toprule
\textbf{Etapa} & \textbf{Operación clave} & \textbf{Producto obtenido} \\
\midrule
Recolección de registros & Guardar observaciones, acciones, estado del entorno y salida del script de evaluación & Permite localizar el punto donde ocurre el fallo \\
Reproducción de trayectorias & Reproducir en orden temporal el comportamiento del Agente y los cambios de la interfaz & Atribución preliminar de la causa del fallo \\
Experimento de control & Sustitución de una sola variable: modelo, prompt, restricción de acciones o versión del script & Evidencia verificable de la corrección \\
\bottomrule
\end{tabular}
\caption{Flujo mínimo viable de depuración para Agentes multimodales}
\end{table}

\begin{warningbox}{Evitar «cambiar muchas cosas a la vez»}
Cuando el modelo, los datos, el formato de las acciones y el script de evaluación cambian al mismo tiempo, resulta casi imposible determinar el origen real de la mejora. El curso enfatiza la iteración ingenieril, que en esencia exige que cada ronda de experimentos sea explicable y reversible.
\end{warningbox}

\subsection{Ejemplos de problemas frecuentes}
\begin{itemize}
    \item \textbf{Desviación en la ubicación de elementos}: los cambios de escala de la interfaz provocan desviaciones de coordenadas; se recomienda introducir ubicación relativa y anclas de objeto;
    \item \textbf{Degradación en tareas largas}: la primera mitad es correcta y la segunda mitad falla; es común en la compresión de contexto y el olvido de la planeación;
    \item \textbf{Manipulación de la recompensa}: el Agente aprende a activar vulnerabilidades del script en lugar de completar la tarea, lo cual exige reforzar la robustez del script de verificación;
    \item \textbf{Desalineación de estado entre aplicaciones}: la ejecución paralela de varias ventanas y tareas en segundo plano provoca inconsistencias de estado, lo cual exige una gestión de sesión más sólida.
\end{itemize}

\begin{importantbox}{Priorización de la depuración}
Corregir primero los fallos «de alta frecuencia y alto costo», y después atender los problemas marginales de baja frecuencia. En un sistema de producción, la ganancia en estabilidad suele valer más que un desempeño óptimo en un solo caso.
\end{importantbox}

\subsection{Resumen de la sección}
La depuración de fallos no es un trabajo accesorio, sino el proceso principal de iteración del Agente. Solo al estructurar los fallos, hacerlos reproducibles y comparables, se forma un verdadero volante de entrenamiento.

